\section{Los límites de rol entre el API Agent y el Web Agent}
\subsection{API Agent: basta una especificación clara para automatizar}
En el segmento de 00:08:35 a 00:09:35, el ponente describe al API Agent como «un conjunto de API o endpoints con nombre que intercambian principalmente información textual»; su conjunto de herramientas puede describirse con precisión mediante una especificación formal, lo cual resulta adecuado para flujos que ya cuentan con una interfaz de servicio.

\begin{knowledgebox}{El beneficio central del API Agent}
Cuando el conjunto de herramientas y los formatos de entrada y salida son claros y verificables, los API Agents pueden mantener una ejecución «de alta precisión y predecible», al mismo tiempo que conservan la reproducibilidad del prompt-engineering.
\end{knowledgebox}

\subsection{Web Agent: requiere percepción y razonamiento sobre la UI}
Los Web Agents se describen como agentes que «hacen clic, llenan formularios y navegan entre páginas dentro del navegador igual que una persona», pero al mismo tiempo cargan con la tarea de entender el objetivo, percibir la situación en el momento y planear a largo plazo (00:32:49--00:34:14).

\begin{warningbox}{El triple desafío del Web Agent}
Tiene que entender la intención del usuario, identificar la UI actual y además planear cómo llegar al resultado final; pero el punto final casi nunca está en la página actual, la observación presente por sí sola no permite determinar si la tarea ya se completó, y un error de ejecución puede provocar que se repitan operaciones o que el proceso se interrumpa.
\end{warningbox}

\subsection{Resumen de la sección}
Los API Agents y los Web Agents no se excluyen entre sí, sino que se complementan: los primeros garantizan estabilidad mediante la especificación, mientras que los segundos logran generalidad mediante la percepción; al diseñar un sistema hay que combinarlos de manera razonable según la tarea objetivo.

